\documentclass[11pt]{article}

\usepackage[margin=1in]{geometry}
\usepackage{amsmath,amssymb,amsthm,mathtools,mathrsfs}
\usepackage{enumitem}
\usepackage{hyperref}
\usepackage{microtype}
\usepackage{bm}
\usepackage{mathtools}

\hypersetup{colorlinks=true,linkcolor=blue,citecolor=blue,urlcolor=blue}

\title{A Proof of the Reimann Hypothesis (WITH GAPS)}
\author{}
\date{}

\theoremstyle{plain}
\newtheorem{theorem}{Theorem}[section]
\newtheorem{proposition}[theorem]{Proposition}
\newtheorem{lemma}[theorem]{Lemma}
\newtheorem{corollary}[theorem]{Corollary}
\newtheorem{claim}[theorem]{Claim}
\newtheorem{gaptheorem}[theorem]{Gap}

\theoremstyle{definition}
\newtheorem{definition}[theorem]{Definition}
\newtheorem{remark}[theorem]{Remark}
\newtheorem{setup}[theorem]{Setup}

\newcommand{\RH}{\mathrm{RH}}
\newcommand{\Om}{\Omega}
\newcommand{\Del}{\Delta}
\newcommand{\Ph}{\Phi}
\newcommand{\Psical}{\Psi}
\newcommand{\zetaSide}{\zeta}
\newcommand{\toy}{\mathrm{toy}}
\newcommand{\val}{\mathrm{val}}
\newcommand{\bg}{\mathrm{bg}}
\newcommand{\sym}{\mathrm{sym}}
\newcommand{\Err}{\mathrm{Err}}
\newcommand{\Tail}{\mathrm{Tail}}
\newcommand{\Tr}{\mathrm{Tr}}
\newcommand{\Ker}{\mathrm{Ker}}
\newcommand{\Span}{\mathrm{span}}
\newcommand{\sgn}{\mathrm{sgn}}
\newcommand{\wt}{\widetilde}
\newcommand{\cK}{c_K}
\newcommand{\cval}{c_{\val}}
\newcommand{\Aval}{A_{\val}}
\newcommand{\Qloc}{Q}
\newcommand{\inner}[2]{\left\langle #1,#2\right\rangle_F}
\newcommand{\norm}[1]{\left\lVert #1\right\rVert}
\newcommand{\abs}[1]{\left\lvert #1\right\rvert}

\begin{document}
\maketitle

\begin{abstract}
This document records a proof strategy for the Riemann Hypothesis based on local jet-normalized blocks of a phase kernel, together with a precise list of the remaining gaps. The text is written in the style of a research proof draft: statements that are already part of the established spine are written as propositions or lemmas with proofs or proof sketches at a level intended to be compatible with a serious mathematical paper; statements that remain unresolved are labeled explicitly as \emph{Gap}. The purpose is to make the architecture of the argument maximally explicit so that the remaining issues can be filled one-by-one.
\end{abstract}

\tableofcontents

\section{Overview}

The argument has three large components.

First, one builds a local matrix statistic from two nearby sample points of the phase kernel and passes from the raw point basis to a point-to-jet basis. In that basis, the local same-point and cross-point blocks have finite limits and can be whitened. This produces a canonical local normalized matrix.

Second, one studies two models for that local matrix: the true zeta/scattering-side matrix and an explicit toy off-line model. On the toy side, the leading deformation is a concrete matrix $M(d)$ depending on a separation parameter $d$, and a certain transverse scalar functional $\Phi_K$ detects it nontrivially for every $d>0$. On the zeta side, after background subtraction and phase-gap subtraction, the local deformation is value-dominated, with curvature controlled by a local estimate of the form
\[
\abs{S''(m)}\ll L(m)S(m)^2.
\]

Third, one uses two scalar functionals in an asymmetric way. A canonical calibration functional, built from the leading zeta-side value-channel matrix, measures the effective size of the local value channel. The functional $\Phi_K$ then serves as a transverse obstruction: it kills the leading zeta value channel but not the toy off-line anomaly. In the end, the zeta-side transverse scalar is quadratically small in the calibrated size, whereas the toy-side transverse scalar is linearly large in that size. This is the intended contradiction.

The only unresolved points are local analytic and quantitative statements. Every such point is isolated as a named \emph{Gap} below.

\section{Set-up and notation}

\begin{setup}[Local midpoint and separation]
Let $m\in \mathbb{R}$ be a local midpoint and let $s\in \mathbb{C}$ be a symmetric separation parameter. Write
\[
t_\pm = m \pm \frac{s}{2}.
\]
Throughout, we use the notation
\[
\Qloc=\Qloc(m)\asymp \log \abs{m}
\]
for the local logarithmic scale. On the background side, one writes
\[
q_\zeta(m)=B_\zeta(m)+S(m),
\]
where $B_\zeta(m)$ is the explicit archimedean/trivial-zero background and $S(m)$ is the positive zero contribution.
\end{setup}

\begin{definition}[Phase kernel]
Let $\Phi$ be a real phase function and define the associated kernel
\[
K_\Phi(x,y):=
\begin{cases}
\dfrac{\sin(\Phi(x)-\Phi(y))}{\pi(x-y)}, & x\neq y,\\[1ex]
\dfrac{\Phi'(x)}{\pi}, & x=y.
\end{cases}
\]
The local zeta/scattering-side phase will be a specific instance of this construction.
\end{definition}

\begin{definition}[Point-to-jet basis]
For a local point spacing $h>0$, define
\[
P_h:=\frac12\begin{pmatrix}1&1\\ -1/h&1/h\end{pmatrix}.
\]
This changes from the raw two-point basis to a value/first-derivative jet basis. The key feature is that local blocks acquire finite limits as $h\to 0$.
\end{definition}

\section{Jet-limit normalization and the local block}

\subsection{The jet-limit same-point block}

Fix a center $T$. Consider the local two-point block $A_h$ built from the pair $T\pm h/2$. Passing to jet coordinates gives
\[
\wt A_h=P_hA_hP_h^\top.
\]
A direct Taylor expansion yields a finite limit.

\begin{proposition}[Jet-limit same-point block]
Let $K=K_\Phi$ be the phase kernel associated to a sufficiently smooth phase $\Phi$. Then
\[
\wt A_h = J(T)+O(h^2),
\]
where
\[
J(T)=
\begin{pmatrix}
2K(T,T) & K_x(T,T)\\[1ex]
K_x(T,T) & \frac12 K_{xy}(T,T)
\end{pmatrix}.
\]
For the phase kernel, writing
\[
q=\Phi',\qquad q'=\Phi'',\qquad q''=\Phi''',
\]
one has
\[
K(T,T)=\frac{q(T)}{\pi},\qquad K_x(T,T)=\frac{q'(T)}{2\pi},\qquad K_{xy}(T,T)=\frac{q''(T)+2q(T)^3}{6\pi}.
\]
Hence
\[
J(T)=\frac1\pi
\begin{pmatrix}
2q(T) & \frac12 q'(T)\\[1ex]
\frac12 q'(T) & \frac1{12}\bigl(q''(T)+2q(T)^3\bigr)
\end{pmatrix}.
\]
\end{proposition}

\begin{proof}
The point-to-jet basis is designed so that the first coordinate records the local value mode and the second the local first-derivative mode. Taylor-expand the kernel entries around $(T,T)$ to second order. The zeroth and first order terms survive in the scaled combination defined by $P_h$, while the residual terms are $O(h^2)$. The stated formulas for $K(T,T)$, $K_x(T,T)$, and $K_{xy}(T,T)$ are obtained by differentiating the explicit phase-kernel formula at the diagonal. Substituting these identities into the matrix definition of $J(T)$ gives the displayed expression.
\end{proof}

\begin{corollary}[Jet Gram positivity criterion]
For the phase kernel,
\[
\det J(T)=\frac{1}{24\pi^2}\Bigl(4q(T)^4+2q(T)q''(T)-3q'(T)^2\Bigr).
\]
In particular, positivity of the local jet Gram matrix is equivalent to the inequality
\[
4q^4+2qq''-3(q')^2\ge 0.
\]
\end{corollary}

\begin{proof}
Compute the determinant directly from the $2\times2$ matrix formula for $J(T)$.
\end{proof}

\subsection{The jet-limit cross block}

Now fix two distinct centers $T_1\neq T_2$, let
\[
q_j=q(T_j),\qquad q_j'=q'(T_j),\qquad q_j''=q''(T_j),\qquad j=1,2,
\]
and write
\[
\Delta:=\Phi(T_1)-\Phi(T_2),\qquad s:=T_1-T_2\neq 0.
\]

\begin{proposition}[Jet-limit cross block]
Let $C_h$ be the local cross block between two two-point samples centered at $T_1$ and $T_2$. Then
\[
\wt C_h=P_hC_hP_h^\top=J_{12}+O(h^2),
\]
where
\[
J_{12}=
\begin{pmatrix}
2K(T_1,T_2) & K_y(T_1,T_2)\\[1ex]
K_x(T_1,T_2) & \frac12 K_{xy}(T_1,T_2)
\end{pmatrix}.
\]
For the phase kernel, write
\[
M_j:=
\begin{pmatrix}
2q_j & \frac12 q_j'\\[1ex]
\frac12 q_j' & \frac1{12}(q_j''+2q_j^3)
\end{pmatrix}
\]
and
\[
N_{12}:=
\begin{pmatrix}
\dfrac{2\sin\Delta}{s} & \dfrac{\sin\Delta-q_2s\cos\Delta}{s^2}\\[2ex]
\dfrac{q_1s\cos\Delta-\sin\Delta}{s^2} & \dfrac{(q_1+q_2)s\cos\Delta+(q_1q_2s^2-2)\sin\Delta}{2s^3}
\end{pmatrix}.
\]
Then
\[
J_j=\frac1\pi M_j,
\qquad
J_{12}=\frac1\pi N_{12}.
\]
Consequently the jet-whitened cross block is
\[
\Omega_\infty(T_1,T_2)=M_1^{-1/2}N_{12}M_2^{-1/2}.
\]
\end{proposition}

\begin{proof}
The proof is parallel to the same-point case. Taylor-expand the phase kernel in the mixed variables around $(T_1,T_2)$ and pass to the point-to-jet basis. The formulas for $K(T_1,T_2)$, $K_x(T_1,T_2)$, $K_y(T_1,T_2)$, and $K_{xy}(T_1,T_2)$ are obtained by direct differentiation of the phase-kernel formula. The common factor $1/\pi$ cancels under whitening.
\end{proof}

\section{The toy anomaly and its transverse obstruction}

The toy off-line model produces a jet-limit anomaly matrix $M(d)$ depending on a separation parameter $d>0$. The explicit formula obtained earlier is
\[
M(d)=
\begin{pmatrix}
\dfrac{d}{(d+i)^3} & \dfrac{i\bigl(2(d+i)^2+1\bigr)}{(d+i)^4}\\[2ex]
\dfrac{2d-i}{(d+i)^4} & \dfrac{4i\bigl((d+i)^2+1\bigr)}{(d+i)^5}
\end{pmatrix}.
\]

\begin{definition}[Transverse obstruction functional]
For a $2\times2$ matrix
\[
X=
\begin{pmatrix}
x_{11}&x_{12}\\x_{21}&x_{22}
\end{pmatrix},
\]
define
\[
\Phi_K(X):=x_{12}+x_{21}.
\]
\end{definition}

\begin{proposition}[Toy transverse obstruction]
For every real $d>0$,
\[
\Phi_K(M(d))
=
\frac{2\bigl(i(d^2-1)-d\bigr)}{(d+i)^4}
\neq 0.
\]
Moreover,
\[
\abs{\Phi_K(M(d))}
=
\frac{2\sqrt{d^4-d^2+1}}{(d^2+1)^2}.
\]
Hence, on every compact interval
\[
D=[d_-,d_+]\subset(0,\infty),
\]
there exists $\cK(D)>0$ such that
\[
\abs{\Phi_K(M(d))}\ge \cK(D)
\qquad(d\in D).
\]
\end{proposition}

\begin{proof}
The formula for $\Phi_K(M(d))$ is obtained by direct addition of the $(1,2)$- and $(2,1)$-entries of $M(d)$. The modulus formula follows by elementary algebra. Since the numerator has nonzero real part $-2d$ for every real $d>0$, the quantity cannot vanish.
\end{proof}

\begin{corollary}[Toy transverse signal]
If the toy deformation has the form
\[
\Delta_\toy(u;d)=u^2M(d)+O(u^4),
\]
then
\[
\Phi_K(\Delta_\toy(u;d))=u^2\Phi_K(M(d))+O(u^4),
\]
and therefore, for $d\in D$,
\[
\abs{\Phi_K(\Delta_\toy(u;d))}\ge \cK(D)u^2-Cu^4.
\]
In particular, for sufficiently small $u$,
\[
\abs{\Phi_K(\Delta_\toy(u;d))}\asymp u^2.
\]
\end{corollary}

\section{Zeta-side value/curvature decomposition}

The zeta-side phase density is decomposed as
\[
q_\zeta(m)=B_\zeta(m)+S(m),
\]
where $B_\zeta$ is explicit and $S(m)$ is the positive zero contribution.

\begin{proposition}[Local curvature estimate]
Let
\[
L(m):=\log\!\Bigl(3+\abs{m}+\frac1{S(m)}\Bigr).
\]
Then
\[
\abs{S''(m)}\ll L(m)S(m)^2.
\]
\end{proposition}

\begin{proof}[Proof sketch]
Write $S$ as a sum of local Poisson-pair profiles. For each pair,
\[
\abs{f_\rho''(m)}\le \frac{24}{\min(\beta_\rho,1-\beta_\rho)}f_\rho(m)^2.
\]
Use the zero-free region to bound the reciprocal of the minimum distance to the critical line by a logarithm in local height, then split the zero sum into near and far parts. The near part contributes $\ll L(m)S(m)^2$, and the far part can be optimized away with a shell-sum argument.
\end{proof}

\section{Finite-$s$ local blocks and whitening}

This section records the abstract finite-$s$ local objects. The exact formulas remain to be written in one canonical notation; this is isolated as Gap~\ref{gap:finite-s-formulas} below.

\begin{definition}[Finite-$s$ local jet blocks]
At the symmetric points
\[
t_\pm=m\pm \frac{s}{2},
\]
let
\[
G_m(s)
\]
denote the same-point $2\times 2$ jet Gram block and let
\[
N_m(s)
\]
denote the mixed left-right $2\times 2$ jet cross block. The whitened local block is
\[
\widehat\Omega_\zeta(s;m)=G_m(s)^{-1/2}N_m(s)G_m(s)^{-1/2},
\]
possibly followed by background and phase-gap subtraction depending on the stage of the argument.
\end{definition}

\begin{remark}
The jet-limit matrices $M_j$ and $N_{12}$ are the limiting avatars of $G_m(s)$ and $N_m(s)$ when one passes to explicit center variables and then to the jet limit. The missing ingredient is the exact finite-$s$ entry formulas in one fixed notation.
\end{remark}

\section{Odd transverse scalar and the all-order expansion}

Define the zeta-side transverse scalar by
\[
H_m(s):=\Phi_K(\widehat\Omega_\zeta(s;m)).
\]
Symmetric placement implies oddness under $s\mapsto -s$.

\begin{proposition}[Oddness of the transverse scalar]
Assume the renormalized local block satisfies the natural jet-basis symmetry
\[
\widehat\Omega_\zeta(-s;m)=J_0\widehat\Omega_\zeta(s;m)J_0,
\qquad
J_0=
\begin{pmatrix}
1&0\\0&-1
\end{pmatrix}.
\]
Then
\[
H_m(-s)=-H_m(s).
\]
\end{proposition}

\begin{proof}
If
\[
X=
\begin{pmatrix}
a&b\\c&d
\end{pmatrix},
\]
then
\[
J_0XJ_0=
\begin{pmatrix}
a&-b\\-c&d
\end{pmatrix},
\]
so
\[
\Phi_K(J_0XJ_0)=-(b+c)=-\Phi_K(X).
\]
Applying this to $X=\widehat\Omega_\zeta(s;m)$ gives the claim.
\end{proof}

The following three gaps are the finite-$s$ analytic package needed to make the odd expansion rigorous.

\begin{gaptheorem}[Exact finite-$s$ formulas and holomorphy]
\label{gap:finite-s-formulas}
Write the exact finite-$s$ entry formulas for $G_m(s)$ and $N_m(s)$ in the actual zeta/scattering normalization, and prove that they are holomorphic for
\[
|s|<\rho_0/\Qloc
\]
with $\rho_0>0$ independent of the admissible midpoint class.
\end{gaptheorem}

\begin{gaptheorem}[Uniform Gram stability and holomorphic whitening]
\label{gap:gram-stability}
Prove the derivative estimate
\[
\sup_{|s|<\rho/\Qloc}\norm{\partial_s G_m(s)}\le C_G\Qloc^2,
\]
and deduce that $G_m(s)$ remains uniformly positive definite on
\[
|s|<\rho_0/\Qloc.
\]
Then prove that
\[
G_m(s)^{-1/2}
\]
is holomorphic there, so the whitened block is holomorphic.
\end{gaptheorem}

\begin{gaptheorem}[Boundary estimate]
\label{gap:boundary}
Prove that on the microscopic boundary
\[
|s|=\rho_0/\Qloc,
\]
one has
\[
\abs{H_m(s)}\ll \frac{L(m)S(m)^2}{\Qloc^3}.
\]
\end{gaptheorem}

Assuming these gaps, Cauchy's theorem gives the odd all-order expansion
\[
H_m(z/\Qloc^2)=\sum_{r\ge 0} c_{2r+1}(m)\frac{z^{2r+1}}{\Qloc^{2r+4}},
\qquad
\abs{c_{2r+1}(m)}\ll L(m)S(m)^2\rho_0^{-(2r+1)}.
\]

\section{$N$-point odd-moment cancellation}

For $N\le \rho_0\Qloc$, define
\[
s_j=\frac{j}{\Qloc^2},\qquad j=1,\dots,N,
\]
and let
\[
a_j^{(N)}
=
(-1)^{j-1}
\frac{2(2N-1)N!(N-1)!j}{(N-j)!(N+j)!}
\]
be the explicit odd-moment-canceling coefficients.

\begin{proposition}[$N$-point cancellation]
These coefficients satisfy
\[
\sum_{j=1}^N a_j^{(N)}=1,
\qquad
\sum_{j=1}^N a_j^{(N)}j^{2r+1}=0
\quad (r=0,\dots,N-2).
\]
Consequently, if the odd analytic expansion of $H_m$ holds, then the combination
\[
\Psi_{\zetaSide}^{(N)}(m):=\sum_{j=1}^N a_j^{(N)}H_m(s_j)
\]
satisfies
\[
\abs{\Psi_{\zetaSide}^{(N)}(m)}
\ll
L(m)S(m)^2\frac{\Theta_N}{\Qloc^{2N+2}},
\]
where $\Theta_N$ is explicit and exponentially tame for $N=\lfloor \alpha \Qloc\rfloor$ with $\alpha>0$ small enough.
\end{proposition}

\begin{proof}[Proof sketch]
The moment-canceling identities are standard for the explicit coefficients. Substituting the odd expansion of $H_m$ shows that the first $N-1$ odd terms vanish, leaving only higher terms. The resulting coefficient sum can be estimated by the exact moment generating function, yielding the displayed bound.
\end{proof}

\section{Corrected endgame: calibration and obstruction}

The recent correction to the endgame is as follows. One does \emph{not} try to make two functionals vanish on the leading zeta value plane. Instead:

\begin{itemize}[nosep]
\item one uses a canonical calibration functional $\Psical$ derived from the actual leading value matrix,
\item and one keeps $\Phi_K$ as the transverse contradiction channel.
\end{itemize}

\subsection{The leading value matrix}

In the symmetric background regime let
\[
T_1=m-\frac{\sigma}{2},\qquad T_2=m+\frac{\sigma}{2},\qquad x:=B_\zeta(m)\sigma.
\]
Let $\Omega_\infty$ be the explicit jet-limit block. Define the leading pure value-channel matrix by
\[
\Aval(m,\sigma):=\partial_\alpha \Omega_\infty\big|_{\alpha=0},
\]
where one perturbs the local value channel by $q_1=q_2=B+\alpha$ and freezes the derivative/curvature data at background level.

The explicit small-$x$ expansion established earlier is
\[
\Aval(x)=\frac1{B}
\begin{pmatrix}
-\dfrac{x^2}{3}+O(x^4) & -\dfrac{\sqrt3 x}{3}+O(x^3)\\[1.2ex]
\dfrac{\sqrt3 x}{3}+O(x^3) & -\dfrac{3x^2}{5}+O(x^4)
\end{pmatrix}.
\]
In particular,
\[
\norm{\Aval(x)}\asymp \frac{x}{B}
\qquad(x\to 0).
\]

Moreover,
\[
\Phi_K(\Aval)=0.
\]
Thus $\Phi_K$ kills the leading zeta value channel.

\subsection{Canonical calibration functional}

Define the canonical dual functional
\[
\Psical_x(X):=rac{\inner{\Aval(x)}{X}}{\norm{\Aval(x)}^2}.
\]
Then
\[
\Psical_x(\Aval)=1.
\]
The small-$x$ pairing with the toy anomaly is
\[
B\,\inner{\Aval(x)}{M(d)}
=
x\,c_1(d)+x^2c_2(d)+O(x^3),
\qquad
c_1(d)=\frac{2\sqrt3 d(id+3)}{3(d-i)^4},
\]
so $c_1(d)\neq 0$ for every real $d>0$.

\begin{proposition}[Stable small-$x$ calibration regime]
Fix a compact interval
\[
D=[d_-,d_+]\subset(0,\infty).
\]
Then there exists $x_0=x_0(D)>0$ and constants $0<c_-(D)\le c_+(D)<\infty$ such that for all
\[
0<x\le x_0,
\qquad
d\in D,
\]
one has
\[
c_-(D)\frac{x}{B}\le \norm{\Aval(x)}\le c_+(D)\frac{x}{B},
\]
and
\[
c_-(D)\frac{x}{B}\le \abs{\inner{\Aval(x)}{M(d)}}\le c_+(D)\frac{x}{B}.
\]
Consequently,
\[
\norm{\Psical_x}\asymp \frac{B}{x},
\qquad
\abs{\Psical_x(M(d))}\asymp \frac{B}{x},
\]
uniformly on this regime.
\end{proposition}

\begin{proof}
For the norm, the off-diagonal entries dominate, as in the explicit small-$x$ expansion; squaring and summing yields
\[
\norm{\Aval(x)}^2=\frac{2x^2}{3B^2}+O\!\left(\frac{x^4}{B^2}\right).
\]
Taking square roots gives the two-sided bound after shrinking $x_0$.

For the pairing, the coefficient $c_1(d)$ is continuous and nonzero on the compact interval $D$, so
\[
m_1(D):=\min_{d\in D}\abs{c_1(d)}>0.
\]
Then
\[
B\abs{\inner{\Aval(x)}{M(d)}}\ge x m_1(D)-C(D)x^2,
\]
and choosing $x_0$ sufficiently small yields the lower bound. The upper bound is immediate from the same expansion.
\end{proof}

\subsection{Zeta-side calibration and transverse obstruction}

The intended zeta-side decomposition is
\[
\Delta_\zeta = S(m)\Aval(x)+R_\zeta,
\]
with
\[
\norm{R_\zeta}\ll L(m)S(m)^2+\Tail+\Err.
\]
Applying $\Psical_x$ gives
\[
\Psical_x(\Delta_\zeta)
=
S(m)+O\!\left(\frac{B}{x}\bigl(L(m)S(m)^2+\Tail+\Err\bigr)\right).
\]
On the toy side,
\[
\Delta_\toy(u;d)=u^2M(d)+O(u^4),
\]
so
\[
\Psical_x(\Delta_\toy(u;d))
=
u^2\Psical_x(M(d))+O(u^4).
\]
Thus $\Psical_x$ calibrates $u^2\leftrightarrow S(m)$.

Meanwhile $\Phi_K$ kills the zeta leading value channel but not the toy anomaly:
\[
\Phi_K(\Delta_\zeta)=O(L(m)S(m)^2)+O(\Tail)+O(\Err),
\]
whereas
\[
\Phi_K(\Delta_\toy(u;d))=u^2\Phi_K(M(d))+O(u^4)\asymp u^2.
\]
This is the contradiction mechanism.

The remaining two gaps isolate the point where the corrected endgame still needs proof.

\begin{gaptheorem}[Canonical calibration theorem]
\label{gap:value-calibration}
In the fixed small-scaled regime
\[
\sigma=\frac{x_0}{B_\zeta(m)}
\]
with $x_0$ as above, prove that the true zeta-side local deformation satisfies
\[
\Delta_\zeta = S(m)\Aval(x_0)+R_\zeta,
\qquad
\norm{R_\zeta}\ll L(m)S(m)^2+\Tail+\Err,
\]
so that
\[
\Psical_{x_0}(\Delta_\zeta)=S(m)+o(S(m)).
\]
\end{gaptheorem}

\begin{remark}
The key step is to prove that the leading value-channel contribution to the true zeta-side deformation is exactly the matrix $\Aval(x_0)$ produced by the finite-dimensional jet-limit linearization.
\end{remark}

\begin{gaptheorem}[Smallness of tail and error terms in the fixed-$x_0$ regime]
\label{gap:smallness-regime}
Prove that in the regime
\[
\sigma=\frac{x_0}{B_\zeta(m)},
\]
one has
\[
B_\zeta(m)L(m)S(m)\to 0,
\qquad
B_\zeta(m)(\Tail+\Err)=o(S(m)).
\]
Equivalently, after calibration, the zeta-side remainder is quadratically small relative to the calibrated value scale.
\end{gaptheorem}

\section{Final conditional theorem}

\begin{theorem}[Conditional proof of $\RH$]
Assume Gap~\ref{gap:finite-s-formulas}, Proposition~\ref{prop:corrected-whitening}, Gap~\ref{gap:boundary}, Gap~\ref{gap:value-calibration}, and Gap~\ref{gap:smallness-regime}, together with Proposition~\ref{prop:tail-curvature}. Then the Riemann Hypothesis holds.
\end{theorem}

\begin{proof}
By Gap~\ref{gap:finite-s-formulas}, Proposition~\ref{prop:corrected-whitening}, and Gap~\ref{gap:boundary}, the odd holomorphic local package is valid and the $N$-point odd-moment cancellation theorem applies. By Gap~\ref{gap:value-calibration}, the canonical calibration functional $\Psical_{x_0}$ identifies the toy size parameter with the zeta value scale:
\[
u^2\asymp S(m).
\]
By construction,
\[
\Phi_K(\Aval)=0,
\]
so the zeta-side transverse signal is only
\[
\Phi_K(\Delta_\zeta)=O(L(m)S(m)^2)+O(\Tail)+O(\Err).
\]
Proposition~\ref{prop:tail-curvature} controls the tail contribution, and Gap~\ref{gap:smallness-regime} controls the remaining quantitative smallness, so this is $o(S(m))$. But on the toy side,
\[
\Phi_K(\Delta_\toy(u;d))\asymp u^2\asymp S(m)
\]
on every compact $d$-interval. Contradiction. Therefore no off-line local deformation can occur. Hence all nontrivial zeros lie on the critical line, i.e. $\RH$ holds.
\end{proof}

\section{Work order}

To complete the proof, the natural order is:
\begin{enumerate}[label=\arabic*.]
\item prove Gap~\ref{gap:finite-s-formulas};
\item prove Gap~\ref{gap:gram-stability};
\item prove Gap~\ref{gap:boundary};
\item prove Gap~\ref{gap:value-calibration};
\item prove Gap~\ref{gap:smallness-regime}.
\end{enumerate}

The first three gaps complete the finite-$s$ local analytic package. Proposition~\ref{prop:tail-curvature} supplies the omitted-zero tail estimate. The last two gaps complete the corrected calibration-plus-obstruction endgame.

\end{document}
